\chapter{Gestión, fiabilidad y soporte logístico}\label{ch:24}
\objetivos{Construir indicadores con denominadores definidos; relacionar recursos y disponibilidad; distinguir una optimización logística de una decisión de aeronavegabilidad.}
\section{Planificar con restricciones}
La gestión organiza demanda de mantenimiento, capacidad, recursos y plazos sin perder las condiciones técnicas y regulatorias. Se identifican tareas obligatorias, vencimientos, tiempos de preparación, repuestos, herramientas, personas y ubicaciones. Agrupar intervenciones puede ahorrar paradas, pero no permite desplazar un requisito fuera de su condición aplicable.

La ruta crítica considera dependencias: desmontar, inspeccionar, evaluar daño, obtener datos, recibir repuesto, montar, probar y certificar. Una demora pequeña en un recurso único puede dominar el tiempo total. La planificación debe conservar incertidumbre y contingencias, no solo una duración ideal.
\section{Indicadores y exposición}
Una tasa $\hat\lambda=n/T$ usa número de eventos $n$ y exposición $T$. Su interpretación depende de definiciones: fallo, retirada programada, retirada no programada, parada en vuelo o demora técnica no son sinónimos. Comparar dos tasas exige poblaciones y periodos equivalentes.

En un modelo de tasa constante, $R(t)=e^{-\lambda t}$ y $MTBF=1/\lambda$. La distribución exponencial no tiene memoria y no describe envejecimiento creciente por el solo hecho de usar horas. Adoptarla es una hipótesis, no una consecuencia automática de haber calculado un promedio.
\begin{equation}R(t)=\exp(-\lambda t),\quad [\lambda]=\mathrm{tiempo}^{-1}.\label{eq:fiabilidadexp}\end{equation}
Con muestras pequeñas, la incertidumbre puede ser grande. Cero eventos en un periodo no demuestra tasa cero. En un modelo de Poisson y un límite unilateral del 95\,\%, cero eventos en exposición $T$ produce la escala $\lambda_U=-\ln(0.05)/T$. Se declara la hipótesis y el tipo de intervalo; no se interpreta como una garantía de seguridad.
\caso{M24. La exposición cambia la comparación}{Una flota ficticia registra 4 eventos en 20\,000 horas y otra 6 en 60\,000. Las tasas son $2\times10^{-4}$ y $1\times10^{-4}\,\mathrm{h^{-1}}$. La segunda tiene más eventos, pero menor tasa observada. Sin intervalos de incertidumbre, equivalencia de población y definición de evento no se concluye que una política sea mejor.
\par\textbf{Desarrollo para comparar tu intento.} Para comparar poblaciones con distinta exposición, calcular eventos/horas en cada una. El recuento absoluto por sí solo omite la duración de observación.

\begin{enumerate}

\item Normalizar los cuatro eventos por su exposición.

\[\fitmath{\mathrm{rateA}=\frac{4}{20000}\approx 0.0002\,\mathrm{1/h}}\]

\item Normalizar los seis eventos por la otra exposición.

\[\fitmath{\mathrm{rateB}=\frac{6}{60000}\approx 0.0001\,\mathrm{1/h}}\]

\item Comparar tasas observadas; el cociente no prueba causalidad ni significación estadística.

\[\fitmath{\mathrm{ratio}=\frac{\frac{6}{60000}}{\frac{4}{20000}}\approx 0.5\,}\]

\end{enumerate}

}
\section{Weibull y tendencias}
La distribución de Weibull se utiliza con fiabilidad $R(t)=\exp[-(t/\eta)^\beta]$. El parámetro de forma modifica la tasa de riesgo: $\beta<1$ permite una tasa decreciente, $\beta=1$ es exponencial y $\beta>1$ creciente dentro del modelo. Un ajuste a datos mezclados puede crear una aparente tendencia que no corresponda a un mecanismo único.

La censura de datos importa: una unidad que no ha fallado al cerrar observación aporta información y no debe tratarse como fallo en esa hora ni excluirse sin más. Los indicadores se acompañan de criterios de recopilación y calidad del dato.
\section{Repuestos y tiempos de espera}
La logística necesita estimar demanda, plazos de reposición, reparabilidad, caducidad y riesgos de ruptura de stock. El stock de seguridad depende de variabilidad y nivel de servicio, no solo de consumo medio. Un modelo de Poisson puede ilustrar demanda discreta, pero su validez debe evaluarse frente a concentración de fallos, campañas y dependencias.

La trazabilidad del repuesto y su elegibilidad no se intercambian por urgencia. Un componente físicamente disponible puede no ser utilizable si no se puede demostrar su estado o compatibilidad. La gestión diferencia material en reparación, servible, en cuarentena y no servible.
\section{Coste y decisiones}
El coste incluye adquisición, intervención, parada, soporte, formación, herramientas y consecuencias de fallos dentro de los criterios de gestión. Un coste esperado puede ayudar a comparar estrategias técnicamente admisibles; no convierte un requisito de seguridad o vida en una opción que se elimina porque resulte cara.

La evaluación de riesgos identifica peligros, probabilidad, severidad y medidas, con incertidumbres y responsabilidades. Una matriz de colores es una herramienta de organización; no sustituye evidencia ni una demostración cuantitativa cuando esta se necesita.
\ilustra{logistica}{Relación entre tarea, recursos, ejecución, repuestos y tiempos de indisponibilidad.}
\section{Cierre del ciclo de gestión}
Se compara plan con ejecución y se investigan desviaciones. Las retiradas repetidas pueden revelar diagnóstico incompleto, dificultades de acceso o datos insuficientes. El seguimiento del cumplimiento y el sistema de gestión de seguridad utilizan esa experiencia con funciones diferenciadas.

Un informe útil identifica acción, responsable, plazo y evidencia de eficacia. Informar una mejoría de un promedio sin comprobar si cambió exposición, definición o población puede ocultar deterioro. La gestión del mantenimiento se apoya en datos técnicamente significativos y decisiones trazables.
\section{Fiabilidad condicional y su interpretación}
Para una Weibull ideal de escala $\eta$ y forma $\beta$, la supervivencia es $R(t)=\exp[-(t/\eta)^\beta]$. Con $\beta=2$ la tasa de riesgo $h(t)=(\beta/\eta)(t/\eta)^{\beta-1}$ crece con la edad del modelo. Si una unidad ha sobrevivido hasta $t_0$, la probabilidad de sobrevivir otras $\Delta t$ es
\[
P(T>t_0+\Delta t\mid T>t_0)=\frac{R(t_0+\Delta t)}{R(t_0)};
\]
no debe calcularse como $R(\Delta t)$ salvo condiciones especiales como la distribución exponencial sin memoria.

Para datos ficticios $\eta=1500\,\mathrm h$, $\beta=2$ y $t_0=1000\,\mathrm h$ resulta $R(1000)\simeq0.64118$; la probabilidad condicional de sobrevivir hasta 1100 h, tras haber llegado a 1000, es $0.91089$. Ningún resultado permite establecer por sí solo un límite de vida autorizado.
\repaso{¿Puede utilizarse un MTBF observado de pocas averías como evidencia suficiente para ampliar vida límite?}{No. Faltan población, dispersión, modelo y base técnica y normativa; MTBF y límite de vida son conceptos distintos.}


\section{Definir el indicador}
Un indicador necesita población, evento, exposición, periodo y reglas. Cinco retiradas no equivalen necesariamente a cinco fallos confirmados. Una tasa dividida por horas de otra población carece de interpretación. Se identifican unidades repetidas, modificaciones y utilización.

Bajo Poisson homogéneo, $\widehat\lambda=n/T$ estima tasa; con supervivencia exponencial, $R(t)=\exp(-\lambda t)$. Proceso de eventos reparables y distribución de vida de una pieza se distinguen. MTBF no es límite de vida ni periodo libre de fallos garantizado. P32 explica por qué cero eventos no prueba tasa cero.
\section{Weibull y observación incompleta}
Con $R(t)=\exp[-(t/\eta)^\beta]$, la tasa instantánea es $h(t)=(\beta/\eta)(t/\eta)^{\beta-1}$. En el modelo crece con edad si $\beta>1$, es constante con $\beta=1$ y decrece con $\beta<1$. La forma estadística no identifica automáticamente mecanismo físico ni una curva de bañera universal.

Las unidades aún operativas aportan censura, no fallos en la fecha de cierre. Ignorarlas puede sesgar la estimación. Censura y pérdida de historial también difieren. Supervivencia condicional usa $R(t_0+\Delta t)/R(t_0)$, generalmente distinta de $R(\Delta t)$ salvo el caso exponencial.
\section{Recursos, repuestos y riesgo}
Planificar exige coincidencia de instalación, personal, útil, datos y repuesto dentro de vencimientos. Agrupar tareas puede reducir paradas y adelantar consumo de intervalos o elevar complejidad. La planificación no desplaza una obligación para mejorar disponibilidad.

Un repuesto físicamente presente puede carecer de elegibilidad o preservación. Suministro, reparación e incertidumbre condicionan stock y parada. Coste de ciclo de vida incluye más que adquisición. Las matrices cualitativas ayudan a priorizar, sin convertir escalas ordinales en probabilidades precisas. Se conserva evidencia y se comprueba eficacia de medidas dentro del sistema pertinente.
\repaso{¿Reducir MTTR excluyendo esperas demuestra mejora operacional?}{Puede haber cambiado solo la definición. Se comparan tiempos de igual alcance y se comprueba si las paradas reales disminuyeron.}

\resumen{Los eventos necesitan exposición y definición. MTBF y tasa no son límites de vida. Exponencial y Weibull son modelos, no etiquetas que prueben un mecanismo. Logística y coste se optimizan entre soluciones admisibles; no eliminan obligaciones o trazabilidad.}
\fuentes{Reglamento 1321/2014, M.A.302(g)--(h), 145.A.47 y .200 \cite{EU1321}; materiales docentes de gestión; FAA--General, factores humanos y registros \cite{FAA-G}. Las distribuciones se presentan como modelos matemáticos explícitos, con comprobaciones en el cuaderno.}
